% 384_FFGFT_Kurzfassung_De_ch.tex
% Johann Pascher, 1. Okt. 2026
% FFGFT in Kurzfassung: Grundlagen, Ergebnisse, Status (Stand nach der Rechenprüfung)
% Überarbeitet am 1. Okt. 2026: Schwerpunkt auf Verhältnisse und natürliche Einheiten
% Erweitert am 2. Okt. 2026: Ikosaeder und phi-Skelett (Dok. 293, 364, 367, 368, 370)
% Neu formatiert am 2. Okt. 2026: Listen, Kästen, farbige Statusmarker; Inhalt unverändert

\providecommand{\BM}{\textbf{[B]}}
\providecommand{\KM}{\textbf{[K]}}
\providecommand{\SM}{\textbf{[S]}}
\providecommand{\XM}{\textbf{[X]}}
\providecommand{\QM}{\textbf{[Q]}}

% Layout (2. Okt. 2026): farbige Statusmarker, Kästen für Kernformeln
\providecommand{\ffst}[2]{\textcolor{#2}{\textbf{[#1]}}}
\renewcommand{\BM}{\ffst{B}{deepgreen}}
\renewcommand{\KM}{\ffst{K}{t0blue!80!black}}
\renewcommand{\SM}{\ffst{S}{t0orange!85!black}}
\renewcommand{\XM}{\ffst{X}{deepred}}
\renewcommand{\QM}{\ffst{Q}{gray!75!black}}
\newtcolorbox{ffkern}{colback=t0blue!4!white,colframe=t0blue!55!black,boxrule=0.5pt,
  arc=2pt,left=4pt,right=4pt,top=0pt,bottom=2pt,before skip=8pt,after skip=8pt,after app={\noindent}}
\newtcolorbox{ffkasten}[1]{colback=gray!4!white,colframe=gray!55!black,boxrule=0.5pt,
  arc=2pt,left=6pt,right=6pt,top=3pt,bottom=3pt,fonttitle=\bfseries\small,title={#1},
  before skip=8pt,after skip=8pt}

\chapter*{FFGFT in Kurzfassung}
\addcontentsline{toc}{chapter}{FFGFT in Kurzfassung}

\begin{abstract}
Die Fundamentale Fraktale Geometrische Feldtheorie (FFGFT) ist in erster Linie eine
verhältnisbasierte Theorie. Sie rechnet in natürlichen Einheiten, $\hbar=c=1$, mit der
Konvention $\alpha=1$ und entsprechend umdefinierter Ladung $e=\sqrt{4\pi}$. In dieser
Grundform gibt es genau einen dimensionslosen Parameter, $\xi=4/30000=1/7500$, und keine
fraktale Korrektur. Teilchen, Kopplungen und Gravitation erscheinen als Eigenschaften eines
kompakten vierdimensionalen Trägers $T^4/\mathbb{Z}_3$; Ausgangspunkt ist die Zeit-Masse-Dualität
$\tilde T\cdot m=1$. Der leitende Gedanke ist Resonanz: Teilchen sind Schwingungsmoden,
Massenverhältnisse sind Periodenverhältnisse, $1/\xi=2^2\cdot3\cdot5^4$ ist ein Punkt im Eulerschen
Tonnetz, und die fraktale Korrektur ist das Komma der Einbettung ins Kontinuum. Die Kernaussagen sind Verhältnisse: $m_\mu/m_e$ und $m_\tau/m_\mu$ folgen
aus $\xi$ auf $0{,}5$ und $1\,\%$, die elektroschwache Skala steht zur Planck-Energie wie
$\xi^4/(5\pi)$, Minimallänge und minimale Zeit stehen zu den Planck-Größen wie $\xi$, die
Neutrino-Massendifferenzen liegen auf $1\,\%$, und die Ladungsquantisierung folgt algebraisch
aus der Galois-Struktur von $\mathrm{GF}(27)^*$. Über das Ikosaeder kommt der goldene Schnitt
hinzu: Die Gewichte der Fünffachdrehung stehen exakt wie $\varphi^8$ und $2\varphi^4$, und
$m_\tau/m_e=(74+\tfrac1{45})\varphi^8$ trifft auf $3\times10^{-5}$, innerhalb der Messunsicherheit. Erst die Übersetzung in SI braucht einen
Messwert als Anker, den Faktor $K_{\text{frak}}=74/75$ und Umrechnungsgrößen wie $C_{\text{conv}}$;
dort entstehen auch die Absolutwerte von $\alpha$ und $G$. Dieses Dokument fasst den Korpus
(Dok.~001--388, A-Serie) auf dem Stand vom 1.~Oktober 2026 zusammen, nach der vollständigen
Nachrechnung. Es nennt für jedes Ergebnis Formel, Zahlenwert, Vergleichswert, Status und tragende
Dokumente und sagt ebenso deutlich, was gesetzt ist, was offen bleibt und was nicht mehr gilt.
Der kosmische Sektor ($H_0$, $\Lambda$) ist bewusst ausgeklammert.
\end{abstract}

% ============================================================
\section*{Statusmarkierungen}
% ============================================================
\begin{center}
\begin{tabular}{@{}lp{11cm}@{}}
\BM & Algebraisch bewiesen (Umformung oder Beweis ohne Zusatzannahme). \\[2pt]
\KM & Numerisch bestätigt (Prüfskript, Vergleich mit Messwerten). \\[2pt]
\SM & Setzung, offen oder spekulativ: motiviert, aber nicht hergeleitet. \\[2pt]
\XM & Widerlegt oder ausgeschlossen. Hinter einem verneinten Satz gilt die Marke der verneinten Behauptung, nicht dem Satz. \\[2pt]
\QM & Konvention oder übernommene Größe. \\
\end{tabular}
\end{center}
Maßgeblich ist der korrigierte Stand: die datierten Hinweise in den Dokumenten
(\glqq Korrigiert am\grqq, \glqq Vermerk vom\grqq, \glqq Entfernt am\grqq) und das
Korrekturregister Dok.~190 bis R147. Vergleichswerte: CODATA 2022, PDG 2024. Alle Zahlen
dieses Dokuments rechnet das Prüfskript \texttt{pruef\_384\_zusammenfassung.py} nach.

% ============================================================
\section{Was FFGFT ist}
% ============================================================

\subsection{Verhältnisse in natürlichen Einheiten}

Der Kern der FFGFT sind Verhältnisse. In natürlichen Einheiten sind Masse, Energie, Temperatur
und inverse Zeit eine Größe in vier Kleidungen (A085), und jede physikalische Aussage lässt sich
als dimensionsloses Verhältnis schreiben: Massenverhältnisse, Skalenverhältnisse zur
Planck-Größe, Mischungswinkel, Ladungsverhältnisse. Solche Verhältnisse hängen an keiner
Einheitenwahl und tragen den Inhalt der Theorie. Absolutwerte in SI sind eine zweite Schicht:
Sie brauchen einen gemessenen Anker und Umrechnungsgrößen, und dort treten Korrekturen auf, die
in den Verhältnissen nicht vorkommen (A080, A130, Dok.~122). Die Kurzfassung folgt dieser
Ordnung: zuerst die Grundform und die Verhältnisse, dann die Übersetzung in SI.

\subsection{Architektur in vier Schichten}

FFGFT ruht auf drei ausdrücklich gesetzten Grundannahmen (R56, Dok.~365): der Zeit-Masse-Dualität
$\tilde T\cdot m=1$, dem kompakten Träger $T^4$ und der Identifikation mit der Ordnung~3,
$\mathbb{Z}_3$. Als vierte, motivierte Strukturwahl kommt das Gitter $D_4$ hinzu; Dok.~380 zeigt,
dass $D_4$ unter den geprüften Vergleichsträgern mit fixpunkthaltiger $\mathbb{Z}_3$-Wirkung die
kleinste Gitterenergie hat \KM, die Wahl bleibt aber eine Setzung \SM. Dazu kommt genau ein
dimensionsloser Parameter, $\xi=4/30000$. Die Feinstrukturkonstante wird als $\alpha=1$ gesetzt,
indem die Ladungseinheit angepasst wird; das ist eine Konvention, kein Ergebnis \QM.
Galois-Körper, Hodge-Theorie und Darstellungstheorie dienen als Konsistenzprüfungen, nicht als
zusätzliche Eingänge (Dok.~365, Schicht~4). Die Raumzeitdimension $3+1$ ist empirische Eingabe
(P34).

\subsection{Resonanz als Grundgedanke: das Eulersche Tonnetz}

Hinter allen Einzelergebnissen steht ein Gedanke: Was existiert, braucht Resonanz. Ein Teilchen
ist ein Schwingungsmodus auf dem kompakten Träger mit ganzen Wicklungszahlen
(Resonanzbedingung, Dok.~189); über $\tilde T\cdot m=1$ wird jede Masse zu einer Frequenz und
jedes Massenverhältnis zu einem Periodenverhältnis. Periodenverhältnisse sind Charaktere
zyklischer Gruppen, und die $\mathbb{Z}_3$-Reduktion macht daraus Arithmetik in
$\mathrm{GF}(3^k)$: Harmonik und Galois-Klassifikation sind nicht analog, sondern zwei
Darstellungen desselben Objekts (Dok.~358, Satz~D \BM). Übergeordnet ist damit die Geometrie, die
Resonanzen besitzt; sie legt alle Verhältnisse fest. Der Messwert ist der Kammerton: Er sagt,
\emph{wo} die Struktur klingt, nicht \emph{wie} (Dok.~310; zum Anker Abschnitt~4).

Das Vorbild ist Eulers Tonnetz (1739), das Gitter der reinen Stimmung mit den Primzahlen
$2$, $3$, $5$ als Achsen. Es trägt die Theorie an drei Stellen:
\begin{itemize}[leftmargin=*,itemsep=2pt]
\item \textbf{$\xi$ ist ein Gitterpunkt.} $1/\xi=7500=2^2\cdot3\cdot5^4$ liegt im 5-Limit-Raum bei
  den Koordinaten $(-2,-1,-4)$ (Dok.~316 \KM). Sieben der neun Yukawa-Vorfaktoren sind
  5-Limit-Verhältnisse wie die Intervalle der reinen Stimmung (Dok.~060); die beiden Ausnahmen,
  die Primzahlen $13$ und $7$, sind in der Galois-Lesart die harmonischen Primen der Tiefe $k=3$
  und $k=6$, alle neun liegen im Raster $\{2,3,5,7,11,13\}$ (Dok.~358, Satz~B \BM).
\item \textbf{Die Massenleiter ist eine logarithmische Resonanzskala.} Wegen $\xi\ll1$ ist die
  Exponentenachse eine Tonleiter: Ein Schritt $\Delta p=1/3$ ist das feste Intervall
  $\xi^{1/3}\approx1/19{,}6$, die \glqq Terz\grqq{} des Teilchenspektrums (Dok.~189).
\item \textbf{Das Komma ist der Einbettungsfehler.} Eulers Spirale aus Quinten und Terzen schließt
  nie, weil $2^x3^y5^z=1$ nur für $x=y=z=0$ gilt; es bleiben Kommas \QM. Im endlichen Körper
  schließt dagegen jeder Quintenzirkel exakt, ein Komma ist dort unmöglich (Dok.~358, Satz~A
  \BM). Das Komma der Musik und die fraktale Korrektur sind deshalb keine Eigenschaft der
  Galois-Struktur, sondern der Einbettung ins Kontinuum; darum erscheint $K_{\text{frak}}$ erst in
  der SI-Übersetzung. Der $\xi$-Zyklus selbst ist eingerollt definiert, mit rationaler
  Schrittweite $1/75$, und schließt nach genau 75 Umläufen (Dok.~315 \BM); die Rotationszahl
  $74/75$ ist derselbe Wert wie $K_{\text{frak}}$.
\end{itemize}
Dieselbe Struktur kehrt in der Kopplung wieder: Unterkritische, kritische und überkritische
Kopplung zweier Kreise entsprechen exakt der $2\times2$-Mischung zweier Zustände (Dok.~328 \BM);
Konsonanz ist Kopplungsgeometrie. Die Grenzen sind ebenso festgehalten: Die harmonische Lesart ist
eine Darstellung, keine zusätzliche Eingabe; ein Tonnetz mit endlichem Limit deckt nur endlich
viele Primzahlen (Dok.~316), und Trefferquoten harmonischer Zuordnungen mit breiten
Toleranzfenstern sind ein Überdeckungseffekt, kein Beleg (Dok.~343).

\subsection{Leitlinien}

Fünf Regeln bestimmen, wie Ergebnisse zu lesen sind:
\begin{enumerate}[leftmargin=*,itemsep=2pt]
\item \textbf{Verhältnisse zählen.} Ein Absolutwert in SI prüft die Umrechnung, nicht die
  Struktur (A130).
\item \textbf{FFGFT kommt mit $\xi$ aus.} Messwerte aus SI und Standardmodell dienen nur der
  Umrechnung und dem Vergleich, und dafür genügt genau ein Messwert als Anker (Dok.~383).
\item \textbf{Ganzzahlige Strukturwahlen sind Setzungen.} $\mathbb{Z}_3$, $D_4$, die
  Wicklungszahlen oder der Exponent~10 im kosmischen Sektor sind motivierte Setzungen, keine
  Kalibrierungen.
\item \textbf{Der kosmische Sektor ist ausgeklammert.} $H_0$ und $\Lambda$ sind bewusst aus der
  Konstantenherleitung ausgeklammert (P39): Das Standardmodell leitet sie ebenso wenig her, und
  keiner der verglichenen Rahmen bestimmt das Verhältnis von Planck-Länge und Hubble-Länge aus
  ersten Prinzipien.
\item \textbf{Erklärung vor Präzision.} Der Ertrag der FFGFT sind Erklärungen: warum Massen in
  festen Verhältnissen stehen, warum die Ladung quantisiert ist, warum es drei Generationen gibt,
  warum Koide $Q=2/3$ gilt, warum die Born-Regel ein geometrischer Quotient ist. Die Zahlen prüfen
  diese Erklärungen; sie ersetzen nicht die Rechnungen des Standardmodells.
\end{enumerate}

\subsection{Was der Korpus heute nicht mehr behauptet}

Die Rechenprüfung hat einige frühere Ansprüche zurückgenommen:
\begin{itemize}[leftmargin=*,itemsep=1pt]
\item FFGFT ist keine Präzisionserweiterung des Standardmodells. Wo dieses rechnet, liefert FFGFT
  dieselben Werte oder benutzt die gemessenen als Anker; eine messbare $\xi$-Korrektur, die genauer
  wäre, gibt es nicht (R142, Dok.~382). Wo es freie Parameter hat, liefert FFGFT Struktur auf
  Prozent- bis Promille-Niveau, nicht auf Messgenauigkeit.
\item FFGFT ist nicht \glqq postulatfrei\grqq: $\tilde T\cdot m=1$, $T^4$, $\mathbb{Z}_3$, $D_4$
  und $\xi$ sind gesetzt.
\item Präzisionstabellen älterer Dokumente, in denen Vorfaktoren an Messwerte angepasst waren,
  sind keine Belege (R138).
\item Die Übereinstimmung des SI-Werts von $\alpha$ ist eine Eigenschaft des gewählten Ankers,
  keine unabhängige Vorhersage (R135).
\item Die G-Präzision auf $0{,}003\,\%$ ist kein eigener Beleg (R141).
\item Ein mit heutiger Hardware messbarer $\xi$-Effekt in Bell-Experimenten oder Quantengattern
  existiert nicht (R142, R143).
\end{itemize}
Wo solche Aussagen in älteren Dokumenten stehen, tragen sie einen datierten Hinweis.

% ============================================================
\section{Die Grundform}
% ============================================================

\subsection{Natürliche Einheiten, \texorpdfstring{$\alpha=1$}{alpha = 1} und die Ladung}

Mit $\hbar=c=\varepsilon_0=1$ (Heaviside-Lorentz) bleibt von der Feinstrukturkonstante die
Beziehung $e^2=4\pi\alpha$ \BM. Die Grundform setzt $\alpha=1$ \QM\ und damit
$e=\sqrt{4\pi}$: Die Ladung wird über die volle Kugeloberfläche gemessen (A267). $\alpha$
verschwindet dabei nicht, es wandert in die Ladungseinheit,
\begin{equation}
  \alpha=\Bigl(\frac{e_{\text{hist}}}{e_{\text{geo}}}\Bigr)^2=\frac{r_e}{\lambda_C}\quad\BM ,
\end{equation}
das Verhältnis von Coulomb- und Compton-Kugel des Elektrons (Dok.~386). $\alpha=1$ heißt
geometrisch: Beide Kugeln fallen zusammen, und mit ihnen der Bohr-Radius. In jeder Rechnung
bleibt $\alpha$ über $e^2=4\pi\alpha$ als Spur erhalten, pro Vertex ein Faktor $\sqrt\alpha$
(A267). Die Brücke zwischen $\alpha$ und $\xi$ läuft über die Energieskala
$E_0=\sqrt{m_em_\mu}$, das geometrische Mittel der beiden leichtesten geladenen Leptonen; in der
Grundform lautet sie
\begin{ffkern}
\begin{equation}
  \xi\,E_0^2=1,\qquad E_0^2=\frac1\xi=7500\quad\BM .
\end{equation}
\end{ffkern}%
Das legt die T0-Energieeinheit fest und lässt keinen Platz für einen Korrekturfaktor.

\subsection{Zeit-Masse-Dualität}

In natürlichen Einheiten sind die intrinsische Zeitskala und die Masse eines Zustands reziprok:
\begin{ffkern}
\begin{equation}
  \tilde T\cdot m=1,\qquad \tilde T=\frac{\hbar}{mc^2}\ \text{(SI)}.
\end{equation}
\end{ffkern}%
Die Relation beschreibt jeden Zustand als Momentaufnahme, nicht als Bewegungsgesetz
(Dok.~365). Eine ortsabhängige Masse im Gravitationspotential ist gleichwertig mit einer
ortsabhängigen Zeit; \glqq Masse konstant, Raum gekrümmt\grqq{} und \glqq Raum flach, Masse
veränderlich\grqq{} sind zwei Darstellungen desselben Sachverhalts (Dok.~078 \KM). Auf dem
Träger gilt die Dualität modenweise exakt als Ortsrelation $\lambda_4\cdot m=2\pi$
(Dok.~314 \KM). Aus der Reziprozität folgt $E\to\infty\Leftrightarrow T\to0$; da die Zeit nach
unten durch $T_0=\xi\,t_P$ beschränkt ist, ist der Horizont ein Extremum und keine
Singularität (Dok.~306).

\subsection{Der Parameter \texorpdfstring{$\xi$}{xi}}

\begin{ffkern}
\begin{equation}
  \xi=\frac{4}{30000}=\frac{1}{7500}=1{,}\overline{3}\times10^{-4}.
\end{equation}
\end{ffkern}%
$\xi$ ist der einzige dimensionslose Parameter. Er hat mehrere konsistente Lesarten, aber keine
erzeugt die Zahl ohne Eingang:
\begin{itemize}[leftmargin=*,itemsep=2pt]
\item \textbf{geometrisch} als Zahl~4 der Gitterrichtungen von $D_4$ über $3\cdot10^4$
  (Dok.~004, 320);
\item \textbf{als Galois-Zahl} über
\begin{equation}
  \frac{(r_\mu/r_e)^2}{\xi}=|\mathrm{GF}(9)^*|^2\cdot5^2\cdot|\mathrm{GF}(27)|=8^2\cdot25\cdot27=43200
\end{equation}
(Dok.~336, 338), was arithmetisch exakt ist \BM, aber an den gesetzten Vorfaktoren $r_i$ hängt;
\item \textbf{spektral} als kleinster Eigenwert des $D_4$-Teiloperators,
  $\xi=\lambda_{\min}(\hat F_{D_4})$ (Dok.~322, 327), der wohldefiniert ist, sofern der
  $D_4$-Sektor im $\mathbb Z_3$-invarianten Sektor liegt \BM, dessen Wert aber eingesetzt und
  nicht hergeleitet ist.
\end{itemize}
Insgesamt ist $\xi$ eine
motivierte Setzung \SM: \glqq Vollständig frei ist $\xi$ damit nicht; vollständig hergeleitet noch
nicht.\grqq{} (Dok.~365).

\begin{ffkasten}{Prüfung über die Higgs-Geometrie}
Als Prüfung dient die Higgs-Geometrie, und auch sie ist ein reines Verhältnis: Mit der
Selbstkopplung $\lambda_h=m_h^2/(2v^2)=0{,}129$ ist
$\xi_{\text{EFT}}=\lambda_h^2v^2/(16\pi^3m_h^2)=m_h^2/(64\pi^3v^2)=1{,}30\times10^{-4}$; der
Ausdruck hängt nur von $m_h/v$ ab und liegt rund $2{,}3\,\%$ unter $4/30000$ (PDG 2024; je nach
Rundung $2{,}2$--$2{,}6\,\%$; Dok.~354, A190, 385) \KM; die Abweichung ist strukturell gedeutet
\SM\ (R59). Das ist eine Konsistenzprüfung auf wenige Prozent, keine exakte Herleitung. Der Weg
geht auf die ursprüngliche Vakuumformel von 2025 mit $\varepsilon_0$ zurück; ausgeschrieben ist sie
$\xi_{\text{EFT}}\cdot\alpha$ und mit $\alpha=1$ derselbe Ausdruck (Dok.~385) \BM.
\end{ffkasten}

\subsection{Skalen als Verhältnisse}

Alle Skalen der Theorie stehen in festen Verhältnissen zueinander. Minimallänge, minimale Zeit und
maximale Energie verhalten sich zu den Planck-Größen wie
\begin{ffkern}
\begin{equation}
  \frac{L_0}{\ell_P}=\frac{T_0}{t_P}=\xi,\qquad \frac{E_{\max}}{E_P}=\frac1\xi=7500 .
\end{equation}
\end{ffkern}%
Der Exponent~1 ist hergeleitet \BM\ (R90): $L_0$ ist die Compton-Reichweite des
Zeitfeld-Mediators mit Masse $m_T=M_{\text{Basis}}/\xi$; die Wahl $M_{\text{Basis}}=m_P$ im
fundamentalen Sektor ist offen \SM. $L_0$ macht die Stufenzahl des fraktalen Operators endlich und
$\hat F$ damit beschränkt \BM\ (Dok.~327). Die elektroschwache Skala steht zur Planck-Energie
und die Leptonmassen über die Leiter (Abschnitt~3) ebenfalls in reinen $\xi$-Verhältnissen
(Dok.~149, 383):
\begin{ffkern}
\begin{equation}
  \frac{v}{E_P}=\frac{\xi^4}{5\pi}=2{,}012\times10^{-17}\ (-0{,}23\,\%),\qquad
  \frac{m_i}{E_P}=\frac{r_i}{5\pi}\,\xi^{p_i+4}\quad\BM .
\end{equation}
\end{ffkern}%
Mit gemessenem $E_P$ ergeben sich die Leptonmassen zu $0{,}504$, $104{,}8$ und $1780{,}9$\,MeV
($-1{,}32$, $-0{,}80$, $+0{,}22\,\%$) \KM. Der Faktor~10 in
$v=E_P/(f^4\cdot\tfrac{\pi}{2}\cdot10)$ mit $f=1/\xi$ ist an den Messwert angepasst. Setzt man statt des
gemessenen $v$ das nackte $v=248{,}93$\,GeV der Leiter ein, verlangt die Formel $10\,K_{\text{frak}}\approx\pi^2$
statt~10; der Faktor~10 wäre dann eine Vermischung von nacktem Leiterwert und SI-Vergleichswert
\SM\ (Dok.~387).

\textit{(Vermerk vom 2.~Okt.~2026 zur Bezeichnung: $f$ ist hier die Grundwindungszahl $f = 1/\xi = 7500$, keine Frequenz. In natürlichen Einheiten ist $f$ nur eine andere Lesart der Dualität $T\cdot m = 1$: $f = t_P/t_0 = \Lambda_{\text{T0}}/E_P$ mit $t_0 = \xi\,t_P$ und $\Lambda_{\text{T0}} = E_P/\xi$, und aus $\xi E_0^2 = 1$ folgt $f = E_0^2$. Als dimensionsloses Verhältnis lässt sich $f$ nicht begründet auf~1 setzen (Dok.~A135). Vgl.\ Dok.~190, R147.)}

\subsection{Zeitstruktur und Rekursion}

Die Zeitstruktur läuft schrittweise:
\begin{equation}
  \xi_{n+1}=\xi_n\,(1-100\,\xi_n),\qquad \xi_0=\tfrac{1}{7500},\qquad
  \xi_n\approx\frac{1}{100(n+75)},\qquad D(k)=\sum_{n<k}100\,\xi_n\approx\ln\!\Bigl(1+\frac{k}{75}\Bigr).
\end{equation}
Im eingefrorenen Fall mit konstantem Defekt $d=1/75$ schließt die Windung nach genau 75
Umläufen; die Rotationszahl ist $74/75=1-100\xi$ \BM. Das ist die geometrische Herkunft des
Werts, der in der SI-Übersetzung als $K_{\text{frak}}$ auftritt (Abschnitt~4). Im laufenden Fall
wird die Zeitwindung zur logarithmischen Spirale mit Selbstähnlichkeitsverhältnis $e$, einer
kontinuierlichen Skaleninvarianz mit Pol bei $n=-75$ (Dok.~295 \BM/\KM). Die Inkremente
$d_n=100\,\xi_n$ fallen von $d_1=1/75$ an monoton; $d_1$ ist das größte, nicht das kleinste
Inkrement, und die untere Zeitschranke ist $T_0=\xi t_P$ (Dok.~306, 307). Dok.~381 zeigt das
Teleskopprodukt $\prod(1-100\xi_k)=\xi_n/\xi_0$ exakt \BM, eine zweite Ordnung
$1/\xi_n\approx100(n+75)+100\ln((n+75)/75)$, die 50- bis etwa 1000-mal genauer ist \KM, und
die Konstante $\psi'(75)=0{,}01342$ zwischen Massen- und Zeitdefekt \KM. Die Zuordnung zu einer
Ein-Schleifen-$\beta$-Funktion hängt von der Schrittkonvention ab \SM. Die Zahl~100 selbst ist
gesetzt, nicht hergeleitet (Vermerk Dok.~133, A040).

\subsection{Geometrie, Hilbertraum, Spektraltheorie}

Der Zustandsraum ist $\mathcal H=L^2(T^4)\otimes\mathbb C^3$; die Übersetzung zwischen
$T^4$ und $\mathcal H$ ist bijektiv und verlustfrei (Dok.~230, 270, 330). Der Torus ist flach:
Die Krümmung verschwindet punktweise über den Metrikabstieg \BM\ (R95). Das Spektrum trennt
die Gitter: $\mathbb Z^4$ beginnt bei $|k|^2=1$ mit Entartung~8, $D_4$ bei $|k|^2=2$ mit 24
(Dok.~314 \KM).

Der fundamentale fraktale Operator $\hat F=\bigoplus_n S_n\circ L_n$ mit $S_n=r_nU_n$ und
$r_n\in(0,1)$ ist beschränkt \BM. Die $\mathbb Z_3$-Paarung macht ihn normal; selbstadjungiert
ist er auf dem $\mathbb Z_3$-invarianten Sektor $P_0$ \BM, mit Defektindizes $(0,0)$ und genau
einer Realisierung. Auf dem ganzen Raum ist nur der Realteil $(\hat F+\hat F^\dagger)/2$
selbstadjungiert; seine Verwendung ist eine Setzung \SM\ (R145, Dok.~327). Die Teilchenmassen
sind Eigenwerte des Massenoperators $\hat M$, nicht von $\hat F$ (Dok.~322, 330).

\subsection{Galois-Struktur}

Die Endlichkörper $\mathrm{GF}(3^n)$ sind die algebraische Sprache der Modenklassifikation
auf $T^4/\mathbb Z_3$. Es gilt $|\mathrm{GF}(9)^*|=8$, $|\mathrm{GF}(27)^*|=26=2\cdot13$,
$|\mathrm{GF}(81)^*|=80=2^4\cdot5$; $\mathrm{GF}(9)$ ist kein Teilkörper von $\mathrm{GF}(27)$,
beide liegen im Kompositum $\mathrm{GF}(729)$. Der Frobenius $x\mapsto x^3$ zerlegt
$\mathrm{GF}(27)^*\cong\mathbb Z_2\times\mathbb Z_{13}$ in zwei Fixpunkte $\{+1,-1\}$ und acht
Dreier-Orbits; auf $\mathbb Z_{13}$ entstehen vier Orbits
\begin{equation}
  O_1=\{1,3,9\},\quad O_2=\{2,5,6\},\quad O_3=\{4,10,12\},\quad O_4=\{7,8,11\}\quad\BM .
\end{equation}
Die 5 ist der einzige primitive Primteiler von $3^4-1$, nicht der kleinste (Dok.~336, 338, 341).
Mit dem IPI-Korrespondenten D.~Matzke ist nur die Algebra gemeinsam bestätigt (Ordnungen in
$G(3)$ und $G(6)$, Dok.~341); die physikalischen Zuordnungen sind FFGFT-seitig.

% ============================================================
\section{Teilchen: die Verhältnisse}
% ============================================================

\subsection{Leptonleiter}

Die geladenen Leptonen folgen der Leiter
\begin{equation}
  \frac{m_i}{v}=r_i\,\xi^{p_i},\qquad r=\Bigl(\tfrac43,\tfrac{16}{5},\tfrac{25}{9}\Bigr),\qquad
  p=\Bigl(\tfrac32,1,\tfrac23\Bigr),\qquad p_g=\tfrac32\bigl(\tfrac23\bigr)^{g-1},
\end{equation}
mit $r_g=n_{\phi,g}^2/n_{\theta,g}$ aus den Wicklungszahlen $(3,2)$, $(5,4)$, $(9,5)$
(Dok.~006, 317, 338) und einer geometrischen Folge der Exponenten mit Startwert $D/2$ und
Verhältnis $2/3$ (Dok.~158, 018). Die Massenverhältnisse enthalten keine Skala:
\begin{ffkern}
\begin{equation}
  \frac{m_\mu}{m_e}=\frac{12}{5}\xi^{-1/2}=120\sqrt3=207{,}85\ (+0{,}52\,\%),\qquad
  \frac{m_\tau}{m_\mu}=\frac{125}{144}\xi^{-1/3}=16{,}99\ (+1{,}03\,\%)\quad\KM .
\end{equation}
\end{ffkern}%
Formel und Verhältnisrechnung sind \BM, die Übereinstimmung \KM; die $r_i$, $p_i$ und
Wicklungszahlen sind Setzungen \SM\ (R138). Die Schrittweite $1/3$ gilt nur von $\mu$ nach $\tau$;
$e$ liegt auf einer eigenen Bahn (Dok.~322). Die Leiter wird ohne $K_{\text{frak}}$ geschrieben;
in Verhältnissen kürzt er sich exakt heraus \BM\ (A040, A100, Dok.~070, 122).

\subsection{Galois-Identitäten, Koide und die vier Wege}

$(m_\mu/m_e)^2_{\text{Leiter}}=43200$ ist arithmetisch exakt; gegen den Messwert bleibt ein Rest
von $-1{,}03\,\%$, eine empirische Diskrepanz \KM, deren Ursache offen ist \SM\ (R113). Die
Koide-Relation, selbst ein reines Verhältnis,
\begin{equation}
  Q=\frac{m_e+m_\mu+m_\tau}{(\sqrt{m_e}+\sqrt{m_\mu}+\sqrt{m_\tau})^2}=\frac23 ,
\end{equation}
gilt für die gemessenen Massen auf $Q-\tfrac23=-0{,}017\,\xi$ (PDG 2024; mit
$m_\tau=1776{,}86$\,MeV $-0{,}046\,\xi$, Dok.~352) \KM; als Vorhersage von $m_\tau$ aus $m_e,m_\mu$
liegt sie $0{,}4\sigma$ neben PDG, nicht entschieden. Aus der Leiter folgt Koide nur
näherungsweise ($+8\,\xi$).

In der $\mathbb Z_3$-Zirkulantenform
$\sqrt{m_k}=c+d\cos(\theta+2\pi k/3)$ gilt $Q=\tfrac23\Leftrightarrow d/c=\sqrt2$ \BM. Die
Matrixelement-Route mit $|A|=\sqrt2/3$ und $\theta=|A|^2=2/9$ (Dok.~367) trifft
$m_\tau/m_e=3477{,}47$ auf $3\times10^{-5}$ \KM, ergibt aber $m_\mu/m_e=206{,}7703$, das
sind 442$\sigma$ neben CODATA; das exakte Paar $(\sqrt2,2/9)$ ist für $m_\mu/m_e$ ausgeschlossen
\XM\ (Dok.~369, R121). Die geometrische Herleitung von $d/c=\sqrt2$ in Dok.~353 trägt nicht \SM.

Dok.~369 vergleicht vier Wege zu den Leptonmassen (Leiter, Koide-Matrixelement, Galois,
$\varphi$-Skelett); Weg~1 und~3 sind für $m_\mu/m_e$ dieselbe Rechnung in zwei Sprachen. Den
$\varphi$-Weg zeigt der folgende Abschnitt.

\subsection{Ikosaeder und goldener Schnitt: das \texorpdfstring{$\varphi$}{phi}-Skelett}

FFGFT ist von Grund auf eine $\mathbb Z_3$-Theorie; der goldene Schnitt kommt über das Ikosaeder
hinzu. Dessen Drehgruppe $A_5$ ist die kleinste endliche Gruppe, in der eine Dreier- und eine
Fünferdrehung zusammen vorkommen, ohne miteinander zu vertauschen (Dok.~370). Die Drei wirkt
dabei als Gitteroperation auf dem kompakten Träger, die Fünf nur als Phase im ausgerollten Raum;
mit dem $D_4$-Gitter ist $A_5$ unverträglich ($5\nmid1152$), im Körper $\mathrm{GF}(81)$ treffen
sich beide \BM\ (Dok.~364, 370).

Die Fünffachdrehung $R_5$ mischt die drei $\mathbb Z_3$-Moden; die
Gewichte $p_j=|\langle v_j\,|\,R_5\,v_e\rangle|^2$ sind
\begin{equation}
  p_0=\frac29,\qquad p_1=\frac{2+3\varphi}{9},\qquad p_2=\frac{5-3\varphi}{9}\quad\BM
\end{equation}
(Dok.~293, A120). Das rationale Gewicht ist die Koide-Phase, $p_0=\theta=2/9$; dasselbe
Matrixelement liefert mit $3\sqrt{p_0}=\sqrt2$ auch die Koide-Amplitude (Dok.~367). Die beiden
$\varphi$-haltigen Gewichte stehen exakt im Verhältnis
\begin{equation}
  \frac{p_1}{p_2}=\varphi^8,\qquad \frac{p_0}{p_2}=2\varphi^4,\qquad
  3\sqrt{p_j}\in\{\sqrt2,\ \varphi^2,\ \varphi^{-2}\}\quad\BM
\end{equation}
(Dok.~368). Das sind reine Zahlen aus der Geometrie, ohne Kosinus und ohne Skala. Die
Massenverhältnisse tragen dieses Skelett; am genauesten trifft
\begin{ffkern}
\begin{equation}
  \frac{m_\tau}{m_e}=\Bigl(74+\frac1{45}\Bigr)\varphi^8=3477{,}469\qquad
  (\text{PDG }3477{,}37\pm0{,}18;\ +3{,}0\times10^{-5},\ 0{,}6\sigma)\quad\KM .
\end{equation}
\end{ffkern}%
\begin{center}
\small
\begin{tabular}{@{}llll@{}}
\toprule
\textbf{Verhältnis} & \textbf{Formel} & \textbf{Wert} & \textbf{gegen PDG} \\
\midrule
$m_\tau/m_e$ & $74\,\varphi^8$ & $3476{,}42$ & $-2{,}7\times10^{-4}$ ($-5{,}3\sigma$) \\
$m_\tau/m_e$ & $(74+\tfrac1{45})\,\varphi^8$ & $3477{,}469$ & $+3{,}0\times10^{-5}$ ($0{,}6\sigma$) \\
$m_\tau/m_e$ & $74\,\varphi^8(1+\tfrac{27}{12}\xi)$ & $3477{,}468$ & $+2{,}9\times10^{-5}$ ($0{,}6\sigma$) \\
$m_\tau/m_e$ & Koide, $\theta=2/9$ & $3477{,}473$ & $+3{,}1\times10^{-5}$ ($0{,}6\sigma$) \\
$m_\mu/m_e$ & $30\,\varphi^4$ & $205{,}62$ & $-0{,}55\,\%$ \\
\bottomrule
\end{tabular}
\end{center}

Auf dem $\varphi$-Weg ist $m_\tau/m_e$ damit so genau wie auf dem Koide-Weg, innerhalb der
Messunsicherheit von $m_\tau$, als reines Verhältnis ohne Anker und ohne $K_{\text{frak}}$. Der
Weg ist algebraisch: Weil $\cos(2/9)$ transzendent ist, kann kein Koide-Verhältnis exakt ein
rationales Vielfaches einer $\varphi$-Potenz sein \BM; die $\varphi$-Form nähert denselben Wert
algebraisch an, die zweite Lage ($\tfrac1{45}$ bzw.\ $\tfrac{27}{12}\xi$) trägt den Rest.

\begin{ffkasten}{Grenzen des $\varphi$-Wegs}
Die Faktoren $74$ und $\tfrac1{45}$ sind beobachtet, nicht aus dem $\mathbb Z_3$/$A_5$-Rahmen
hergeleitet \SM\ (Dok.~368). Die~$74$ ist dieselbe Zahl wie im Zähler der Rotationszahl $74/75$
der Rekursion (Dok.~381); ihr Primfaktor~37 kommt nachweislich nicht aus dem Ikosaeder, weil
$\mathrm{GF}(81)$ kein Teilkörper von $\mathrm{GF}(3^{18})$ ist \BM\ (Dok.~370). Für $m_\mu/m_e$
gibt es keinen $\varphi$-Weg auf diesem Niveau: $30\varphi^4$ liegt $0{,}55\,\%$ daneben, die
Leiter $0{,}52\,\%$, und das exakte Koide-Paar ist dort ausgeschlossen (oben). Das
$\varphi$-Skelett ist nicht die frühere $\varphi$-Massenleiter mit $m\propto\varphi^{\text{Generation}}$
und $m_\mu/m_e\approx\varphi^{10}$, die widerlegt ist \XM\ (Dok.~172, 281).
\end{ffkasten}

\subsection{Neutrinos}

Dok.~340 leitet Hierarchie und Mischung aus den vier Frobenius-Orbits ab. Die Massenverhältnisse
sind fest,
\begin{equation}
  m_1:m_2:m_3=1:\sqrt{\tfrac{14}{3}}:11,\qquad
  \frac{\Delta m^2_{\text{atm}}}{\Delta m^2_{\text{sol}}}=\frac{120}{11/3}=32{,}7\ (-1{,}4\,\%)\quad\KM ,
\end{equation}
in normaler Hierarchie. Für Absolutwerte braucht es eine Basismasse; angesetzt ist
$m_\nu=\tfrac{\xi^2}{2}m_e=4{,}54$\,meV \SM. Damit folgen $m_2=9{,}81$ und $m_3=49{,}96$\,meV,
$\sum m=64{,}3$\,meV, $\Delta m^2_{\text{atm}}=120\,m_\nu^2=2{,}476\times10^{-3}$\,eV$^2$ ($-1{,}0\,\%$
gegen $2{,}500\times10^{-3}$, Dok.~340) und $\Delta m^2_{\text{sol}}=\tfrac{11}{3}m_\nu^2=7{,}565\times10^{-5}$\,eV$^2$
($+0{,}46\,\%$) \KM. Der Reaktorwinkel ergibt $\sin^2\theta_{13}=(25\xi)^{2/3}=0{,}0223$ ($+1{,}4\,\%$)
\KM. Nach R108 liegen $\nu_1$ und $\nu_2$ algebraisch in der selbstinversen Majorana-Schicht und
$\nu_3$ ist Dirac; diese Zuordnung von $\nu_3$ zum Orbit $O_4$ steht in Spannung zu Dok.~346, das
$O_4$ den geladenen Leptonen zuordnet, und Dok.~344 führt alle Neutrinos als Dirac \SM; dass die
algebraische Selbstinversion physikalisch $\nu=\bar\nu$ bedeutet, ist eine weitergehende Annahme
\SM. Daraus folgt für den neutrinolosen Doppelbetazerfall $m_{ee}\approx6{,}0$\,meV (Phasen null),
destruktiv $0{,}13$\,meV, höchstens $14{,}4$\,meV, und exakt null, falls alle Neutrinos Dirac sind
(Dok.~344). Ältere Neutrinospektren (Dok.~006, 007) sind ausgeschlossen \XM\ (R109; Vermerk
Dok.~007). Das $A_5$-Achterspektrum gibt $\sin^2\theta_{23}\in\{4/9,5/9\}$; $5/9=0{,}556$ liegt
nahe am Messwert der normalen Hierarchie, maximale Mischung $1/2$ ist damit ausgeschlossen
(Dok.~372) \KM. $\theta_{12}$ ist nur der Größenordnung nach erfasst, die CP-Phase nicht \SM. Der
Venting-Exponent $p=5/24$ folgt aus der Zerlegung $\rho_4\!\downarrow_{A_4}=\rho_3\oplus\rho_1$ mit
$|A_5:A_4|=5$ \BM\ (Dok.~340).

\subsection{Anomale magnetische Momente}

Tragend ist die Verhältnislinie \KM\ (Dok.~018, 158): Mit Windungsexponenten $5/3$ und $4/3$ gilt
$\Delta a(\tau{-}\mu)/\Delta a(\mu{-}e)=f^{1/3}-1=18{,}57$, und verankert an den gemessenen
$a_e$, $a_\mu$ folgt
\begin{equation}
  a_\tau=a_e+\frac{144}{125}\frac{m_\tau}{m_\mu}(a_\mu-a_e)=1{,}2811\times10^{-3},
\end{equation}
$1{,}04\times10^{-4}$ über dem SM-Wert; heutige Schranken sind rund 85-mal weiter, die
Entscheidung steht aus \SM. Ein fester FFGFT-Zusatz $\Delta a_\mu=251\times10^{-11}$ als
Erklärung der früheren Myon-Anomalie ist nach dem Stand 2025 überholt; er läge $3{,}4\sigma$
über der heutigen Differenz \XM\ (Dok.~382).

\subsection{Elektroschwacher Sektor, Higgs, Top}

Auch hier sind die Aussagen Verhältnisse zur $Z$-Masse. On-shell gilt
$\sin^2\theta_W=1-M_W^2/M_Z^2=2/9$ als Element des $A_5$-Achterspektrums (Dok.~372) \KM; daraus
folgt $M_W/M_Z=\sqrt{7/9}$, $M_W=80{,}42$\,GeV, $3{,}8\sigma$ vom Weltmittel 2024, eine
ausgewiesene, falsifizierbare Spannung. Im $\overline{\text{MS}}$-Schema liefert der
Renormierungslauf von $3/8$ aus $\sin^2\theta_W(M_Z)=0{,}2308$ ($-0{,}19\,\%$) mit zwei externen
Eingaben (Dok.~323) \KM. Die Verhältnisse $m_h/M_Z=\tfrac{11}{8}$ ($m_h=125{,}38$\,GeV,
$+0{,}15\,\%$) und $m_t/M_Z=(\tfrac{11}{8})^2$ ($m_t=172{,}40$\,GeV, $-0{,}10\,\%$) sowie die
Spurregel $M_W^2+M_Z^2+m_h^2=v^2/2$ ($+0{,}45\,\%$) sind Kandidaten \KM, die Herkunft der Faktoren
ist Lesart \SM\ (Dok.~373). Die Kette \glqq Symmetriebrechung topologisch erzwungen\grqq{} ist an
zwei Stellen offen \SM\ (Dok.~355). Frühere Higgs-Formeln ($m_h=v\xi^{1/4}$, Dok.~041) sind
widerlegt \XM\ (R122).

\subsection{Quarks}

Die Quarks folgen derselben Leiterform $m_q/v=r_q\xi^{p_q}$ mit topologisch begründeten
Exponenten \KM, außer für $c$ und $t$, deren Exponenten Setzungen sind \SM\ (R138); die
Vorfaktoren $r_q$ sind aus Messwerten bestimmt, ihre Herleitung ist offen \SM\ (R129, Dok.~373).
Die Reproduktion ist daher eine Kodierung, kein Beleg; die Tabelle gibt die Werte mit
$v=246$\,GeV:
\begin{center}
\small
\begin{tabular}{@{}lcccc@{}}
\toprule
\textbf{Quark} & $r_q$ & $p_q$ & \textbf{Leiter} ($v=246$\,GeV) & \textbf{PDG 2024} \\
\midrule
$u$ & $6$ & $3/2$ & $2{,}27$\,MeV & $2{,}16$\,MeV \\
$d$ & $25/2$ & $3/2$ & $4{,}73$\,MeV & $4{,}70$\,MeV \\
$s$ & $26/9$ & $1$ & $94{,}8$\,MeV & $93{,}5$\,MeV \\
$c$ & $2$ & $2/3$ & $1{,}284$\,GeV & $1{,}273$\,GeV \\
$b$ & $3/2$ & $1/2$ & $4{,}26$\,GeV & $4{,}18$\,GeV \\
$t$ & $1/28$ & $-1/3$ & $172{,}0$\,GeV & $172{,}6$\,GeV \\
\bottomrule
\end{tabular}
\end{center}
Leichte Quarkmassen sind schemaabhängig; $\tilde T\cdot m=1$ spricht von Polmassen.
$\Lambda_{\text{QCD}}$ folgt nicht aus der $\xi$-Leiter \XM\ (R122): $v\,\xi^{1/3}=12{,}6$\,GeV,
nicht 200\,MeV. Der CKM-Parameter $\lambda=\xi^{1/6}=0{,}2260$ liegt $0{,}46\,\%$ neben dem
Messwert \KM; der Exponent $1/6$ ist eine motivierte Setzung \SM\ (Vermerk Dok.~336). Der Quark-
und Hadronsektor insgesamt ist eine offene Brücke (Dok.~318, R76).

\subsection{Ladung, Hyperladung, Generationen, Farbe, Photon}

Die Ladungsquantisierung folgt algebraisch \BM\ (Dok.~346, R111): Die QR-Orbits $O_1$, $O_3$
mod~13 sind elektrisch neutral (Neutrinos, Antineutrinos), die NQR-Orbits $O_4$ und $O_2$ geladen
(geladene Leptonen, Up-Quarks); zusammen mit $N\bmod3$ folgt
$Q\in\{0,\pm\tfrac13,\pm\tfrac23,\pm1\}$, und \glqq neutral und farbig\grqq{} ist ausgeschlossen.
Die Ladungsverhältnisse sind damit ebenfalls Verhältnisse ganzer Zahlen. Die Down-Quarks haben
keinen eigenen Orbit. Gell-Mann--Nishijima ist nicht vollständig geschlossen: $I_3=\pm\tfrac12$
folgt \BM, die Hyperladung der geladenen Leptonen folgt aus dem QR/NQR-Bit allein nicht und ist
offen \SM\ (R146, Dok.~347). Die drei Elemente eines Frobenius-Orbits geben drei Generationen
\BM\ (Dok.~348); die frühere Vertauschung der Generationen~2 und~3 ist ein Ordnungsartefakt.
$\mathfrak{su}(3)$ entsteht aus den $\mathbb Z_3$-Eigensektoren \BM; $SU(3)$ ist dabei die
kleinste, nicht die einzige passende Gruppe (Dok.~321). Photon und Gluonen sind masselos \BM;
dass die Masselosigkeit Spin~1 bedeutet, folgt daraus nicht und ist offen \SM\ (Dok.~339).

% ============================================================
\section{Die Übersetzung in SI}
% ============================================================

\subsection{Ein Anker}

Weil alle Skalen in festen $\xi$-Verhältnissen stehen (Abschnitt~2), genügt für die
Übersetzung in SI ein einziger Messwert (Dok.~383). Mit $m_e$ als Anker folgen
$E_P=1{,}237\times10^{19}$\,GeV ($+1{,}34\,\%$), $G=6{,}50\times10^{-11}$ ($-2{,}6\,\%$),
$\ell_P=1{,}595\times10^{-35}$\,m und $L_0=\xi\ell_P$ \KM. Die Streuung je nach Wahl des Ankers
ist der bekannte Rest der Leiter, kein zweiter Anker:
\begin{center}
\small
\begin{tabular}{@{}lccc@{}}
\toprule
\textbf{Anker} & $E_P$ & $G$ & $\ell_P$, $L_0$ \\
\midrule
$m_e$ & $+1{,}34\,\%$ & $-2{,}62\,\%$ & $-1{,}32\,\%$ \\
$m_\mu$ & $+0{,}81\,\%$ & $-1{,}60\,\%$ & $-0{,}80\,\%$ \\
$m_\tau$ & $-0{,}22\,\%$ & $+0{,}45\,\%$ & $+0{,}22\,\%$ \\
$v$ & $+0{,}23\,\%$ & $-0{,}46\,\%$ & $-0{,}23\,\%$ \\
\bottomrule
\end{tabular}
\end{center}
\noindent Damit liegen $L_0=2{,}155\times10^{-39}$\,m, $T_0=7{,}19\times10^{-48}$\,s und
$E_{\max}=9{,}16\times10^{22}$\,GeV fest. Mit $v=246$\,GeV liegen die absoluten Leptonmassen bei
$-1{,}18$, $-0{,}66$ und $+0{,}37\,\%$.

\subsection{Elektroschwache Skala}

$v$ ist in der Massenrechnung eine Zwischengröße. Je nach Weg ergibt sich \KM:
\begin{center}
\small
\begin{tabular}{@{}lc@{}}
\toprule
\textbf{Weg} & $v$ \\
\midrule
Galois-$m_e$ aus Dok.~338 (R104) & $248{,}3$\,GeV \\
gemessenes $m_e$ & $248{,}9$\,GeV \\
gemessenes $m_e$ und $K_{\text{frak}}$ & $245{,}6$\,GeV \\
$\xi^4E_P/(5\pi)$ & $245{,}65$\,GeV \\
\bottomrule
\end{tabular}
\end{center}
Der Vergleichswert
$v=246{,}22$\,GeV ist eine SM-Definition aus $G_F$ \QM\ (R112); konventionsfrei wäre nur ein
Vergleich über die Myon-Lebensdauer $\tau_\mu$, der offen ist \SM.

\subsection{Fraktale Korrektur}

\begin{ffkern}
\begin{equation}
  K_{\text{frak}}=1-100\,\xi=\frac{74}{75}=0{,}98\overline{6}.
\end{equation}
\end{ffkern}%
Der Faktor tritt nur in der SI-Übersetzung auf: In Verhältnissen kürzt er sich exakt heraus
\BM\ (A040, Dok.~070, 122), und in der Grundform $\xi E_0^2=1$ hat er keinen Platz. Sein Wert
stammt aus der Geometrie, als Rotationszahl der eingefrorenen Rekursion (Abschnitt~2). Bestätigt
ist er \KM\ (R135): Der Vergleich der beiden Wege zu $E_0$ mit dem gemessenen $\alpha$ verlangt
$K=\xi\,m_em_\mu\,\alpha^{-1}=0{,}98650$, also $n=101{,}2$ statt $100$ in $1-n\xi$. Dieser Faktor
ist genau das Quadrat der nackten Bezugsenergie in MeV (Abschnitt~4.4, Dok.~386, A130). Der Rest
von $1{,}7\times10^{-4}$ ist durch Messunsicherheiten nicht erklärt und offen \SM. Die
\emph{Form} des Faktors (additiv $1-100\xi$ oder multiplikativ $(1-\xi)^{100}$) ist wieder offen
\SM\ (R139): Die verfügbaren Zeugen trennen die beiden Formen nicht schärfer als ihr Abstand von
$8{,}8\times10^{-5}$. Gleichwertig lässt sich $K_{\text{frak}}$ als Potenzform
$(D/3)^{D/2}=74/75$ mit $D_f^{\text{eff}}=2{,}97303$ schreiben \KM\ (Dok.~381); die räumliche
Box-Dimension $D_f=3-\xi=2{,}999867$ ist eine Setzung \SM, die aus der 100er-Rekursion nicht
folgt (Vermerk Dok.~322).

\subsection{\texorpdfstring{$E_0$}{E0} und der SI-Wert von \texorpdfstring{$\alpha$}{alpha}}

In SI ist $E_0$ das geometrische Mittel der beiden leichten Leptonmassen, nackt und fraktal
korrigiert (R72):
\begin{ffkern}
\begin{equation}
  E_0^{\text{nackt}}=\sqrt{m_em_\mu}=7{,}348\ \text{MeV},\qquad
  E_0=\sqrt{\frac{m_em_\mu}{K_{\text{frak}}}}=7{,}397\ \text{MeV}.
\end{equation}
\end{ffkern}%
Mit $\alpha_{\text{SI}}=\xi\,(E_0/1\,\text{MeV})^2$ folgt ohne Korrektur $\alpha^{-1}=7500/E_0^2=138{,}9$,
mit Korrektur $137{,}06$, $1{,}7\times10^{-4}$ neben dem Messwert. Die Galois-Form
$\alpha^{-1}=3700/27=137{,}037$ ($+7{,}6$\,ppm, Dok.~338, R102) folgt aus
$(E_0^{\text{nackt}})^2=m_em_\mu\approx54\ \text{MeV}^2=|\mathrm{GF}(3)^*|\cdot|\mathrm{GF}(27)|\ \text{MeV}^2$
zusammen mit $K_{\text{frak}}=74/75$. Beide Übereinstimmungen sind Eigenschaften des Ankers
$E_0$, keine unabhängigen Vorhersagen \KM\ (R135).

Die SI-Form verlangt die Bezugsenergie
$E_0\sqrt{\xi/\alpha}=0{,}99992$\,MeV (nackt $0{,}99323$\,MeV, das ist $\sqrt{K_{\text{frak}}}$\,MeV).
Das MeV ist keine geometrische Größe, sondern eine Einheit der SI-Kette \QM: Sein Zahlenfaktor
ist der Zahlenwert der historischen Ladung $e$, also derselben Einheit, in der $\alpha$ steckt.
Legt man die Korrektur auf diese Seite, die Seite der Coulomb-Kugel, bleiben die Massen und die
Compton-Kugel nackt, und es ergibt sich $1$\,MeV bis auf $8\times10^{-5}$; diese Zuordnung ist eine
Setzung \SM, der Rest bleibt offen (Dok.~386).
Dritte Wege zu $E_0$ und geschlossene $\alpha$-Formeln aus $\xi$ allein sind zurückgenommen \XM.

\subsection{Gravitationskonstante}

\begin{equation}
  G_{\text{SI}}=\frac{\xi^2}{4m_e}\,C_{\text{conv}}\,K_{\text{frak}},\qquad
  C_{\text{conv}}=7{,}783\times10^{-3},\quad K_{\text{frak}}=0{,}986 .
\end{equation}
$G$ ist keine unabhängige Konstante: Form und Größenordnung folgen aus $\xi$ und dem einen Anker
\KM\ (Dok.~012, 180, 383). $C_{\text{conv}}$ ist die SI-Umrechnung der Ein-Anker-Kette; ihr
Korpuswert enthält zusätzlich den kleinen Rest dieser Kette (Kettenwert $7{,}58\times10^{-3}$).
Die Übereinstimmung mit CODATA auf $0{,}003\,\%$ ergibt sich erst mit diesem Rest und ist kein
eigener Präzisionsbeleg (R141); sie gilt mit dem gerundeten Korpuswert $K_{\text{frak}}=0{,}986$
der G-Formel, mit $74/75$ wären es $+0{,}07\,\%$. Den Rest strukturell zu erklären ist offen \SM.
Aus $\tilde T\cdot m=1$ folgt $m\cdot G=\xi^2/4$; eine Massensingularität mit endlichem $G$ ist
ausgeschlossen \BM\ (Dok.~350), die physikalische Gültigkeit im Regime Schwarzer Löcher ist
offen \SM.

% ============================================================
\section{Gravitation und Kosmos}
% ============================================================

\subsection{Schwarze Löcher}

Setzt man Compton- und Schwarzschild-Radius gleich, folgt rein geometrisch
\begin{equation}
  M_{\text{coll}}=\frac{m_P}{\sqrt2},\qquad r_s=\sqrt2\,\ell_P,\qquad
  S_{\text{BH}}(M_{\text{coll}})=2\pi\ \text{nat}=\frac{2\pi}{\ln2}\ \text{bit}=9{,}065\ \text{bit}\quad\BM
\end{equation}
(Dok.~325, 329). $S_{\text{BH}}$ entspricht Matzkes Länge in Planck-Einheiten; Matzkes Schwelle ist
um $\sqrt2$ kleiner, $n_{\text{thr}}=6{,}41$ (R144). Als universelle Bitzahl ist $n_{\text{thr}}$
widerlegt \XM\ (R88). Das Verhältnis von Sektor- zu thermischer Information ist massenunabhängig
$\ln3\approx1{,}099$ \BM. Die Hawking-Temperatur entspricht der Bit-Energie $\hbar c/L$ an der
Horizontskala \KM; emittierbar sind nur Moden mit Trialität~0, Confinement und Hawking-Selektion
sind ein Prinzip \BM\ (Dok.~325). Graukörper-Faktoren und Rückreaktion sind offen.

\subsection{Casimir und CMB}

Die Casimir-Energiedichte zwischen idealen Platten ist $-\pi^2\hbar c/(720\,d^4)$; die Form mit
240 ist der Druck. Bei der Länge $L_\xi$ ist das Verhältnis zur CMB-Dichte
$\pi^2/(720\xi)=102{,}8$ eine Identität, weil $L_\xi$ selbst aus der CMB-Dichte und $\xi$
bestimmt wird; ein Messbeleg ist es nicht (R136). Einen Faktor $4/3$ liefert die
QFT-Casimir-Energie nicht \XM\ (Dok.~371). Die Temperatur
$k_BT_{\text{CMB}}=\tfrac{16}{9}\xi(1-\tfrac{275}{4}\xi)$\,eV trifft nur in der Einheit eV und
mit nicht hergeleitetem Koeffizienten \SM\ (R137); das eV bringt die Elektronenmasse als
zweiten Anker mit. In der Grundform hängt $T_{\text{CMB}}/m_e$ an einer einzigen Zahl. Der
Kandidat $T_{\text{CMB}}/m_e=(8\pi)^{1/4}\xi^{5/2}$ trifft den FIRAS-Wert auf $0{,}12\sigma$,
ergibt mit der $H_0$-Form $T_{\text{CMB}}^4=16\,H_0m_e^3$ und legt $L_\xi$ ohne die CMB fest;
hergeleitet ist der Vorfaktor nicht \SM\ (Dok.~388). Der Faktor~3 vor der Peak-Folge
$\{1,6,14,26\}$ ist eine Setzung; die Folge selbst ist aus den Daten abgelesen (R139).

\subsection{Kosmischer Sektor: bewusst ausgeklammert}

Die Strukturform $H_0/\nu_e=\pi^2\xi^{10}=1{,}75\times10^{-38}$ ergibt $H_0=66{,}8$\,km/s/Mpc;
der Exponent~10 ist eine motivierte Setzung (P20), die Aussage daher bedingt und keine
Vorhersage. Ebenso bedingt sind $\Lambda^*=1/R_H^2$ und die Zerlegung der $10^{123}$-Frage in
eine Strukturfrage (Dok.~312). Ausgeschlossen sind die $\Lambda$-Formeln älterer Dokumente,
$H_0=\xi c$, $z=\xi x$ und $z\approx\xi\ln(d/\ell_P)$ \XM\ (R137; Vermerk Dok.~156). Die Rotverschiebung
$1+z=e^{\xi x}$ ist achromatisch \KM, die Länge in $x$ ist aber nicht aus $\xi$ ausgewiesen; die
Supernova-Zeitdehnung ist $1+z$ für alle $z$ \BM\ (R128). Supernovae, BAO und CMB-Temperatur sind
zwischen Expansion und intrinsischer Zeitentfaltung entartet; die Entscheidung ist theoretisch,
nicht empirisch (Dok.~267).

\subsection{Galaxien}

Die Beschleunigungsskala $a_0=c^2\xi^{10}/(4\bar\lambda_e)=1{,}03\times10^{-10}$\,m/s$^2$ liegt
$14\,\%$ unter dem empirischen Wert und erbt die Setzung~P20. Die gasdominierte Zwerggalaxie
DDO\,154 trifft auf $0{,}2\,\%$ \KM; für Spiralgalaxien bleiben Abweichungen von rund $19\,\%$
offen. \glqq Bullet Cluster bestanden\grqq{} ist nicht belegt: Mit baryonischen Massen liegt das
Maximum am Gas (Vermerk Dok.~308). Die Lichtablenkung reproduziert den GR-Wert $1{,}75''$ \KM.

% ============================================================
\section{Quantenmechanik}
% ============================================================

FFGFT ist in der Quantenmechanik deterministisch. Die Rechenregeln und Zahlenwerte der
Standard-QM bleiben, ihre Deutung ändert sich an vier Stellen:
\begin{itemize}[leftmargin=*,itemsep=2pt]
\item \textbf{Messung ist ein Polübergang, kein Kollaps.} Die Wechselwirkung mit dem Apparat
  treibt den Zustandspunkt zu einem der Pole; welcher, legt die Apparatgröße $\lambda$ auf
  derselben Achse fest, $A(z,\lambda)=\mathrm{sgn}(z-\lambda)$. Jede Einzelmessung ist damit
  festgelegt, auch wenn $\lambda$ unbekannt ist (Dok.~230, 175) \KM.
\item \textbf{Die Born-Regel ist ein geometrischer Quotient.} $(1+z)/2=|\alpha|^2$ ist der Anteil
  der Achse $[-1,1]$, dessen $\lambda$-Werte zum Pol $+1$ führen \KM; nach Archimedes ist das
  zugleich ein Flächenanteil der Bloch-Kugel \BM. Als Häufigkeit erscheint er erst in
  Messreihen mit unkontrolliertem $\lambda$. Der Zufall ist epistemisch, nicht ontologisch
  (Dok.~365, 367, 370); der nächste Verwandte in der Literatur ist das Hidden-Measurement-Modell
  von Aerts \QM.
\item \textbf{Gatter sind geometrische Drehungen.} Die drei Pauli-Matrizen entsprechen drei
  Rotationen auf dem Bloch-Zylinder; CNOT ist eine geometrische Konditionierung, bei der die
  Drehung des einen Zylinderpunkts vom $z$-Wert des anderen abhängt. Da $\{H,T,\mathrm{CNOT}\}$
  universell ist, lässt sich jedes Gatter deterministisch und ohne Wiederholung nachbilden, bis
  zur Rauschgrenze (Dok.~230, 175) \KM.
\item \textbf{Rauschen ist ungeordnete Phase, und die Grenze gilt für beide Auswertungen.} Ein
  Spektrum mit festen Amplituden und unaufgelösten Phasen sieht aus wie weißes Rauschen, obwohl
  jede Komponente festliegt, und mit bekannten Phasen lässt es sich exakt zurückrechnen \KM;
  Vakuumrauschen kann so eine Überlagerung deterministischer
  Harmonischer sein \SM. Der Ablauf ist deterministisch und nachrechenbar bis zur Rauschgrenze:
  In der Rückkopplung wächst der Phasenfehler mit jedem Schritt und erreicht sie nach etwa
  $\pi/(2\varepsilon)$ Schritten (Dok.~183) \SM. Darüber lassen sich Ausgänge nicht mehr
  unterscheiden. Wiederholte Messungen braucht die deterministische Lesart nicht, denn jeder
  Einzelausgang ist festgelegt. Die probabilistische Auswertung kommt aber auch nicht weiter:
  Die Mittelung über viele Läufe stößt an denselben Boden, dessen natürliche Skala $\xi$ ist
  (Dok.~343, Satz~G$'$) \SM.
\end{itemize}
Bell und Kochen--Specker schließen nur lokale bzw.\ nicht-kontextuelle verborgene Variablen aus \QM.
Die FFGFT-Messung ist kontextuell, weil die Festlegung in der Wechselwirkung liegt, und für
mehrere Qubits ist die Korrelation eine Eigenschaft der geschlossenen Geometrie, keine lokale
Größe im $\mathbb{R}^3$ (Dok.~230, R126). Operativ liefern FFGFT und Standard-QM dieselben
Vorhersagen bis auf eine $\xi$-Korrektur.

Die Bell-Verletzung bleibt: Die fraktale Dämpfung verkleinert $S$ nur um einen Betrag der
Ordnung $\xi$, je nach Form $\Delta S/S\approx10^{-5}$ bis $10^{-4}$, und stellt keine lokale
Schranke her (R142). Die einzige dokumentierte Hardware-Messung (IBM, 28.~Mai 2026, Dok.~147,
2048 Shots je Lauf):
\begin{center}
\small
\begin{tabular}{@{}lcccc@{}}
\toprule
\textbf{Gerät} & \textbf{Läufe} & $\langle S\rangle$ & \textbf{Standardfehler} & \textbf{Anteil an }$2\sqrt2$ \\
\midrule
ibm\_kingston & 15 & $2{,}7396$ & $0{,}0071$ & $96{,}9\,\%$ \\
ibm\_marrakesh & 10 & $2{,}7016$ & $0{,}0138$ & $95{,}5\,\%$ \\
gepoolt & 25 & $2{,}7244$ & $0{,}0078$ & $96{,}3\,\%$ \\
\bottomrule
\end{tabular}
\end{center}
\noindent Die Bell-Verletzung ist bestätigt \KM, die Lücke zu $2\sqrt2$ (bei ibm\_kingston
$0{,}088$, gepoolt $0{,}104$) ist Dekohärenz, 160- bis mehrere tausendmal größer als der
$\xi$-Effekt, und zwei baugleiche Chips unterscheiden sich um $0{,}038$. Der $\xi$-Effekt ist mit
heutiger Hardware nicht auflösbar \SM.

Früher angegebene CHSH-Werte aus angeblichen Läufen von
2025 und ein daran angepasster $\xi$-Wert sind unbelegt und entfernt (R143). Die Abweichung des
$\pi$-Gatters um $100\xi=1{,}33\,\%$ wird bei der Pulskalibrierung echter Qubits mit eingestellt
und ist nicht direkt messbar (R142).

\textbf{Shor in der deterministischen Lesart.} Shors Algorithmus zerfällt in die Aufbereitung,
eine modulare Exponentiation mit $O(n^3)$ Gattern und fast den gesamten Kosten (Anteil
$1-1/(n+1)$, bei 2048~Bit $99{,}95\,\%$ \KM), und die
Quanten-Fourier-Transformation mit $O(n^2)$ Gattern, eine Wellenoperation, die die Periode als
Resonanz herausholt (Dok.~176). Theoretisch faktorisiert er in $O((\log N)^3)$ \KM; ein
allgemeiner exponentieller Quantenvorteil ist dagegen nicht belegt (Dok.~176). Der Vorteil sitzt allein
in der Periodenfindung. Aufbereitung und Fehlerkorrektur kosten auf Quantenhardware mehr als
klassisch: die Aufbereitung um einen festen Faktor, die Fehlerkorrektur um einen Aufwand, der
nach dem Schwellentheorem nur langsam wächst. Wirksam würde der Vorteil erst bei Zahlen, für
die heute die Hardware fehlt (RSA-2048: $8{,}6\times10^9$ Gatter, $4{,}1\times10^6$ Qubits \KM); im
erreichbaren Bereich faktorisiert ein klassischer Rechner ebenso schnell. Als Apparat ist der
Quantenrechner ein Analogrechner mit quantenmechanischer Präzision: Ein Qubit ist ein
Zweiniveau-Resonator, jedes Gatter ein resonanter Antrieb, und die Fehlerkorrektur mittelt über
Syndrom-Runden. Diese Mittelung stößt an denselben Auflösungsboden wie jede Auswertung; dass
sie ihn bei Skalierung unterschreitet, behauptet das Schwellentheorem, gezeigt ist es nicht
\SM. Ob netto ein Vorteil bleibt, entscheidet damit allein die Fehlerkorrektur. Ist $\xi$ der
Boden, liegt die Grenze eines resonanzbasierten Rechners bei $N_\mathrm{max}\approx1/\xi=7500$,
rund 13~Bit, und der Zeitvorteil verschwindet strukturell \SM\ (Dok.~343, Satz~G$'$).
Für die Rechenoperationen selbst gibt es Alternativen: Die Fourier-Transformation läuft optisch
oder in jedem Resonator (Dok.~176), analoge Matrixfaltung und Resonator-Einrasten rechnen mit
Resonanzen statt mit Gattern (Dok.~186), ebenso Resonanzrechner wie die BAW-Ising-Maschine
(Dok.~362). Sie unterliegen demselben Auflösungsboden. Die Rechnungen der Standard-QM bleiben
dabei verwendbar; FFGFT deutet sie als Resonanzvorgänge. Zurückgezogen ist nur die frühere
Behauptung, ein eigenes T0- bzw.\ $\xi$-Faktorisierungsverfahren liefere von Shor
unterscheidbare Ergebnisse (R65); $\xi$ als Linienbreite der Phasenschätzung ist für den
vorregistrierten Aufbau widerlegt, die Linie bleibt offen (P44).

% ============================================================
\section{Prüfbare Aussagen}
% ============================================================

\begin{center}
\small
\rowcolors{2}{gray!7}{white}
\begin{longtable}{@{}p{3.4cm}p{3.6cm}p{3.6cm}p{3.2cm}@{}}
\toprule
\textbf{Größe} & \textbf{FFGFT} & \textbf{Messwert / Grenze} & \textbf{Status} \\
\midrule
\endhead
\rowcolor{t0blue!12}\multicolumn{4}{@{}l}{\textbf{Verhältnisse (Grundform)}} \\
$m_\mu/m_e$ & $207{,}85$ & $206{,}77$ ($+0{,}52\,\%$) & \KM, $r_i$ \SM \\
$m_\tau/m_\mu$ & $16{,}99$ & $16{,}82$ ($+1{,}03\,\%$) & \KM \\
Koide $Q$ & $2/3$ & $-0{,}017\,\xi$ & \KM \\
$m_\tau$ aus Koide & $1776{,}97$\,MeV & $1776{,}93(9)$, $0{,}4\sigma$ & offen (Fußnote~1) \\
$m_\tau/m_e$ ($\varphi$-Skelett) & $(74+\tfrac1{45})\varphi^8=3477{,}47$ & $3477{,}37(18)$, $0{,}6\sigma$ & $\varphi^8$ \BM, Faktoren \SM \\
$v/E_P$ & $\xi^4/(5\pi)$ & $-0{,}23\,\%$ & \KM, Faktor 10 angepasst \\
$\xi$ aus $m_h/v$ & $1{,}30\times10^{-4}$ & $4/30000$ ($-2{,}3\,\%$) & Prüfung \KM \\
$\Delta m^2_{\text{atm}}/\Delta m^2_{\text{sol}}$ & $32{,}7$ & $33{,}2$ ($-1{,}4\,\%$) & \KM \\
$\sin^2\theta_{13}$ & $0{,}0223$ & $0{,}0220$ ($+1{,}4\,\%$) & \KM \\
$\sin^2\theta_{23}$ & $5/9$ & normale Hierarchie & \KM \\
$a_\tau$ & $1{,}2811\times10^{-3}$ & Schranke 85-mal weiter & offen (Fußnote~2) \\
$M_W/M_Z$ & $\sqrt{7/9}$, $80{,}42$\,GeV & $3{,}8\sigma$ (Weltmittel) & Spannung \\
$m_h/M_Z$ & $11/8$, $125{,}38$\,GeV & $125{,}20$ ($+0{,}15\,\%$) & Kandidat \KM \\
$m_t/M_Z$ & $(11/8)^2$, $172{,}40$\,GeV & $172{,}57$ ($-0{,}10\,\%$) & Kandidat \KM \\
$\lambda_{\text{CKM}}$ & $\xi^{1/6}=0{,}2260$ & $0{,}2250$ ($+0{,}46\,\%$) & \KM, Exponent \SM \\
CHSH-Abweichung & $\Delta S/S\sim10^{-5}$--$10^{-4}$ & Hardware-Rauschen $\sim10^{-2}$ & nicht auflösbar \\
$\pi$-Gatter & $1{,}33\,\%$ & in Pulskalibrierung & nicht messbar \\[3pt]
\rowcolor{t0blue!12}\multicolumn{4}{@{}l}{\textbf{Absolutwerte (SI-Übersetzung)}} \\
$\Delta m^2_{\text{atm}}$ & $2{,}476\times10^{-3}$\,eV$^2$ & $2{,}500\times10^{-3}$ ($-1{,}0\,\%$) & \KM \\
$\Delta m^2_{\text{sol}}$ & $7{,}565\times10^{-5}$\,eV$^2$ & $+0{,}46\,\%$ & \KM \\
$m_{ee}$ ($0\nu\beta\beta$) & $6{,}0$ / $0{,}13$ / $0$\,meV & nEXO 5--20\,meV & \SM, entscheidet nEXO \\
$\sum m_\nu$ & $64{,}3$\,meV & kosmologisch, modellabhängig & \KM \\
$\alpha^{-1}$ & $137{,}06$ & $137{,}036$ ($+1{,}7\times10^{-4}$) & Eigenschaft des Ankers \\
$\tau_\mu$ ohne $G_F$, $v$ & -- & MuLan & offen \SM \\
\bottomrule
\end{longtable}
\end{center}
{\small
\noindent\textbf{Fußnote 1.} Koide mit den gemessenen $m_e$ und $m_\mu$ ergibt $m_\tau=1776{,}97$\,MeV,
$0{,}4\sigma$ neben dem PDG-Mittel $1776{,}93(9)$\,MeV \KM. Entschieden ist der Test damit nicht:
Die beiden Messmethoden, Schwellenscan (BESIII, $1776{,}91$\,MeV) und Pseudomasse (Belle~II,
$1777{,}09$\,MeV, mit $m_{\nu_\tau}=0$), streuen um $0{,}18$\,MeV, das Doppelte der angegebenen
Unsicherheit. Erst wenn die Methoden übereinstimmen, prüft der Wert Koide (R113, Dok.~351, 352).

\noindent\textbf{Fußnote 2.} $a_\tau$ ist noch nicht gemessen; es gibt nur ein Ausschlussintervall.
Die beste Schranke ist CMS 2024 mit $-0{,}0042<a_\tau<0{,}0046$ (95\,\% CL), also $0{,}0088$ breit.
Der FFGFT-Wert liegt $1{,}04\times10^{-4}$ über dem Standardmodell-Wert $1{,}1772\times10^{-3}$; das
Intervall ist rund 85-mal breiter als dieser Unterschied, beide Werte liegen darin \KM. Entscheiden
lässt sich die Frage erst mit einer etwa hundertfach genaueren Messung, etwa an Belle~II
(Dok.~382).\par}

% ============================================================
\section{Was nicht mehr gilt}
% ============================================================

Die folgenden Aussagen älterer Dokumente sind widerlegt oder zurückgenommen \XM; die
Dokumente tragen dazu datierte Hinweise.
\begin{description}[leftmargin=0pt,itemsep=3pt,font=\normalfont\bfseries]
\item[Grundlagen:] $\xi$ als zwingend hergeleitet,
ein Casimir-Faktor $4/3$, die Zahl~100 als hergeleitet, ein minimales
Zeitinkrement $d_1$, die Selbstadjungiertheit von $\hat F$ auf dem ganzen Raum, ein Körperturm
$\mathrm{GF}(3)\subset\mathrm{GF}(9)\subset\mathrm{GF}(27)$.
\item[Teilchen:] Präzisionstabellen mit angepassten Vorfaktoren als Belege, die
$\varphi$-Massenleiter mit $m\propto\varphi^{\text{Generation}}$ (nicht das $\varphi$-Skelett), das exakte Paar $(\sqrt2,2/9)$ für $m_\mu/m_e$, $\Lambda_{\text{QCD}}$ aus
der $\xi$-Leiter, frühere Higgs-Formeln, die Neutrinospektren von Dok.~006 und 007, ein fester
Myon-Zusatz $\Delta a_\mu=251\times10^{-11}$, Gell-Mann--Nishijima als vollständig geschlossen.
\item[Gravitation und Kosmos:] die G-Präzision als eigener Beleg, $n_{\text{thr}}$ als universelle Zahl,
das Casimir-CMB-Verhältnis als Messbeleg, die $\Lambda$-Formeln, $H_0=\xi c$ und $z=\xi x$,
\glqq Bullet Cluster bestanden\grqq.
\item[Quantenmechanik:] lokaler Realismus oder
$S\le2+\varepsilon$, eine Hardware-Bestätigung des $\xi$-Effekts, die CHSH-Werte 2,8275 und
2,8278, ein FFGFT-eigenes Faktorisierungsverfahren.
\end{description}

% ============================================================
\section{Offene Brücken}
% ============================================================

Dok.~190 führt die offenen Punkte in einer Brückenliste (Stand R147). Zusammengefasst:
\begin{enumerate}[leftmargin=*,itemsep=1pt]
\item \textbf{Grundform:} die Herkunft der Zahl~100 und des angepassten Faktors~10 in $v/E_P$ (Dok.~149, 387);
  $M_{\text{Basis}}=m_P$ (R90); die Selbstadjungiertheit von $\hat F$ auf dem ganzen Raum (R145);
  ein adversarieller $D_4$-Spezifitätstest (R95); die Spektraldimension $d_s=1{,}86$ gegen
  $\approx2$.
\item \textbf{SI-Übersetzung:} der Rest $1{,}7\times10^{-4}$ zwischen dem von $\alpha$ verlangten
  $K_{\text{frak}}$ und $1-100\xi$ und die Form des Faktors (R135, R139); der Rest von
  $8\times10^{-5}$ bei der Bezugsenergie $1$\,MeV (Dok.~386); der Rest der Ein-Anker-Kette bei $G$ (R141).
\item \textbf{Teilchen:} $\tau_\mu$ direkt aus FFGFT als einziger konventionsfreier Test des schwachen
  Sektors (R112); die Ursache der fallenden Leptonreste (R113, R121); die Koide-Amplitude
  $d/c=\sqrt2$ aus der Geometrie; die Faktoren $74$ und $\tfrac1{45}$ des $\varphi$-Skeletts und
  die Frage, ob ikosaedrische Phase und Galois-Ordnung eine gemeinsame Wurzel haben (Dok.~368,
  370); die Hyperladung der geladenen Leptonen (R146); die
  Quark-Vorfaktoren und der Hadronsektor (R76, R129); $K_{\text{had}}$ (R123); zwei-Schleifen-Lauf
  von $\alpha_s$; $m_{\text{Pl}}$ und $\alpha_{\text{em}}(M_Z)$ aus $\xi$; Spin~1 des Photons.
\item \textbf{Gravitation und Kosmos:} Graukörper und Rückreaktion (R85); eine physikalische
  Casimir-Signatur von $L_\xi$ (R136); $T_{\text{CMB}}$ einheitenunabhängig (R137); die
  CMB-Peak-Selektion vorwärts und die $C_\ell$-Projektion (P29--P31); der Exponent $41/4$ bzw.\
  10 vorwärts (P20); leichte Elemente ohne heiße Frühphase; $\Omega_{\text{DM}}$ modellneutral.
\item \textbf{Methodik:} $\xi$ als Linienbreite (P44); Sprachbereinigung zu $H_0$ (P16, P34).
\end{enumerate}

% ============================================================
\section{Statusbilanz}
% ============================================================

\begin{center}
\small
\rowcolors{2}{white}{gray!7}
\begin{longtable}{@{}p{1.2cm}p{13cm}@{}}
\toprule
\textbf{Status} & \textbf{Ergebnisse (Auswahl)} \\
\midrule
\endhead
\BM & $e^2=4\pi\alpha$ und $\alpha=r_e/\lambda_C$; $\xi E_0^2=1$ in der Grundform (Dok.~386);
  Ladungsquantisierung aus dem Legendre-Symbol (Dok.~346); $SU(2)_L$, $U(1)_Y$ und
  $\sin^2\theta_W|_{\text{GUT}}=3/8$ (R80); $\mathfrak{su}(3)$ aus $\mathbb Z_3$-Sektoren
  (Dok.~321); drei Generationen als Frobenius-Orbit (Dok.~348); Selbstadjungiertheit von $\hat F$
  auf $P_0$ (R145); Exponent in $L_0=\xi\ell_P$ (R90); $m_i/E_P=\tfrac{r_i}{5\pi}\xi^{p_i+4}$
  (Dok.~383); Rotationszahl $74/75$ und Teleskopprodukt der Rekursion (Dok.~381);
  $A_5$-Gewichte mit $p_0=\theta=2/9$, $p_1/p_2=\varphi^8$, $p_0/p_2=2\varphi^4$ (Dok.~293, 367, 368);
  $S_{\text{BH}}(M_{\text{coll}})=2\pi$\,nat und $\ln3$ (R144); Supernova-Zeitdehnung $1+z$
  (R128). \\[3pt]
\KM & Leptonverhältnisse und Koide; $m_\tau/m_e$ aus dem $\varphi$-Skelett auf $3\times10^{-5}$; Neutrino-Verhältnisse, Massendifferenzen und $\theta_{13}$;
  $a_\tau$-Verhältnisse; $v/E_P$; Higgs-Prüfung von $\xi$ auf rund $2{,}3\,\%$; Higgs- und
  Top-Kandidaten; $\lambda_{\text{CKM}}$; Wert $K_{\text{frak}}=74/75$ und sein Auftreten nur in
  der SI-Übersetzung; Form und Größenordnung von $G$; Bell-Verletzung auf Hardware. \\[3pt]
\SM & $\tilde T\cdot m=1$, $T^4$, $\mathbb Z_3$, $D_4$, $\xi$; $r_i$, $p_i$ und Wicklungszahlen;
  die Faktoren $74$ und $\tfrac1{45}$ im $\varphi$-Skelett;
  die Zahl~100 und die Form von $K_{\text{frak}}$; Faktor~10 in $v/E_P$ ($10\,K_{\text{frak}}\approx\pi^2$ mit nacktem $v$); $M_{\text{Basis}}=m_P$;
  Korrektur auf der Einheitenseite (Dok.~386); Basismasse $m_\nu$; Quark-Vorfaktoren; Hyperladung der geladenen
  Leptonen; Spin~1 des Photons; Exponent~10 im kosmischen Sektor; CMB-Faktor~3. \\[3pt]
\QM & $\hbar=c=1$ und $\alpha=1$ mit angepasstem $e$; $v=246$\,GeV als SM-Definition; der Anker
  $m_\mu/m_e$ ist QED-abhängig; Maxwell-Dynamik übernommen; $C_{\text{conv}}$ und das
  MeV als SI-Umrechnung. \\[3pt]
\XM & siehe Abschnitt \glqq Was nicht mehr gilt\grqq. \\
\bottomrule
\end{longtable}
\end{center}

% ============================================================
\section{Methodik und Buchführung}
% ============================================================

Der Korpus führt jede Aussage mit einem Statusmarker. Neuere Dokumente haben ein Prüfskript mit
Assertions, viele davon Nulltests. Das Register Dok.~190 bucht Korrekturen an älteren
Dokumenten und offene Brücken; seit R133 werden die Quelldokumente selbst korrigiert, jede Stelle
mit datiertem Hinweis und dem vorherigen Wert. Ausgeschlossene Vorhersagen werden nicht
umgeschrieben, sondern vermerkt. Vom 29.~September bis 1.~Oktober 2026 wurde der ganze Korpus
nachgerechnet: rund 235 Dokumente tragen seitdem einen Korrekturkasten, die A-Serie trägt
Hinweise im Text. Die Folgen für die Einstufung stehen in R135 bis R146. Leitplanken:
\begin{itemize}[leftmargin=*,itemsep=1pt]
\item Ein Vergleichswert ist kein Eingang.
\item Eine Form $X=\xi^N$ allein sagt nichts (P35).
\item Vorhersagen werden mit festgelegten Freiheitsgraden vorab registriert (P44).
\item Jede Vorhersage nennt ihren Anker (R113).
\item Ein Absolutwert in SI prüft die Umrechnung, ein Verhältnis die Struktur (A130).
\end{itemize}

% ============================================================
\section{Leseführer}
% ============================================================

\begin{center}
\small
\rowcolors{2}{gray!7}{white}
\begin{tabular}{@{}p{3.2cm}p{11cm}@{}}
\toprule
\textbf{Thema} & \textbf{Dokumente} \\
\midrule
Fundament & 365 (vier Schichten), 190 (Register), 180/189 (Kette $\xi\to L_0$), 383 (ein Anker) \\
Einheiten, $\alpha$ & A080/A085 (natürliche Einheiten), A130/A267 ($\alpha=1$), 122 (Verhältnis und Absolutwert), 386 (wo $\alpha$ steckt), 385 (Higgs-Vakuum-Weg), 387 (Abweichungen) \\
Massen & 006 (Leiter), 336/338 (Galois), 351/352 (PDG, Leptonreste), 367/369 (Matrixelement, vier Wege) \\
Ikosaeder, $\varphi$ & 291/292/293 (Ort und Herkunft von $2/9$), 364 ($A_5$ und $D_4$), 368 ($\varphi$-Skelett), 370 (warum das Ikosaeder) \\
Elektroschwach & 323 (Weinberg-Lauf), 372/373 (2/9, Higgs, Yukawa), 355 (Symmetriebrechung) \\
Ladung, Farbe & 321 ($SU(3)$), 346/347 (Ladung, Gell-Mann--Nishijima), 348 (Generationen) \\
Neutrinos, $g-2$ & 340, 343/344 (Majorana/Dirac), 158/382 ($a_\tau$, Stand 2025) \\
Gravitation & 012/013 ($G$, SI), 325/329 (Schwarze Löcher), 350, 388 (CMB-Temperatur in der Grundform) \\
Struktur & 314 (Gitter), 322/327/330 (Spektraltheorie), 380 ($D_4$), 381 (Rekursion) \\
Zeit & 295 (Log-Spirale), 306/307 (Reziprozität, Zeit im Zustandsraum) \\
Quantenmechanik & 230/231 (Hilbertraum, Messung), 175 (Qubit-Zustandsräume, deterministisches Auslesen), 183 (Rauschen als Phase), 023 (Bell), 147 (Hardware), 176 (Shor: Aufbereitung und Transformation), 343 (Quantenrechner als Analogrechner) \\
Resonanz, Harmonik & 060 (Musik-Parallelen), 189 (Resonanzbedingung), 310 (Resonanzgeometrie), 315 (Euler-Spirale), 316 (Tonnetz), 328 (Kopplungsregime), 358 (Harmonik und Galois) \\
Abgrenzung & 309 (Skalenanker), 267 (Entartung), 377 (Schnittstellen), 378 (Begutachtung) \\
\bottomrule
\end{tabular}
\end{center}

% ============================================================
\section{Fazit}
% ============================================================

\begin{ffkern}
FFGFT ist eine Theorie der Verhältnisse und liefert Erklärungen, keine größere Genauigkeit als
das Standardmodell. In natürlichen Einheiten mit $\alpha=1$ ersetzt sie die
freien Parameter des Leptonsektors durch eine Struktur mit einem Parameter $\xi$ und gesetzten
ganzen Zahlen, stellt Leptonmassen, elektroschwache Skala, Minimallänge und Planck-Skala in feste
$\xi$-Verhältnisse und gewinnt die Ladungsquantisierung algebraisch. Ihr Grundgedanke ist Resonanz:
Teilchen sind Moden des Trägers, $\xi$ ist ein Punkt im Eulerschen Tonnetz, und die fraktale
Korrektur ist das Komma der Einbettung. Die Verhältnisse treffen auf
Prozent bis Promille; über das Ikosaeder trägt $m_\tau/m_e$ das $\varphi$-Skelett und trifft auf
$3\times10^{-5}$. Die Übersetzung in SI braucht einen Anker, die fraktale Korrektur und
Umrechnungsgrößen; dort entstehen die Absolutwerte von $\alpha$ und $G$, und dort liegt ein
bekannter Rest. Offen sind die Herkunft der Zahl~100, die Herkunft des angepassten Faktors~10, die Form der fraktalen
Korrektur und ihr Rest gegenüber dem MeV, die Hyperladung der geladenen Leptonen, der
Hadronsektor und alles, was am kosmischen Sektor hängt. Entscheidend prüfbar sind in den nächsten
Jahren $m_{ee}$ (nEXO), $a_\tau$ (Belle~II), die Neutrinomassensumme und $M_W$. Der Korpus führt
Setzungen als Setzungen und Widerlegungen als Widerlegungen; diese Kurzfassung gibt diesen Stand
wieder.
\end{ffkern}

% ============================================================
\section{Verwendung von KI-Werkzeugen}
% ============================================================

Dieses Dokument ist unter Verwendung KI-gestützter Werkzeuge entstanden. Die
Fragestellungen, die Setzungen und alle inhaltlichen Entscheidungen stammen vom
Autor; die Werkzeuge formulieren aus, rechnen nach und erstellen Prüfskripte.
Die Verantwortung für Inhalt und Fehler liegt vollständig beim Autor.

\vspace{1em}
\noindent\textit{Prüfskript:} \texttt{2/python/Dok384\_Skripte/pruef\_384\_zusammenfassung.py}
--- alle Zahlen dieses Dokuments nachgerechnet.
